% Shared by the architecture white paper and the modem implementation reference.
\providecommand{\modemcite}[1]{}

\subsection{Audio ingress, continuity, and candidate detection}
The optional CPAL input backend requires native 48,000 samples/s and a selected
zero-based channel; it accepts float32, signed 16-bit, or unsigned 16-bit
formats and converts to normalized mono float samples. It selects one channel
rather than averaging potentially phase-opposed channels. No implicit
resampling is implemented. The callback accumulates 1,024-sample blocks and
uses a nonblocking bounded queue of 64 blocks. This represents approximately
1.365 s of mono input; it is a queue capacity, not a permitted DSP or keying
latency. Queue overflow, a malformed interleaved input block, or device failure
latches a terminal capture fault. The receiver must restart on a new continuous
stream instead of concealing dropped samples.

DSP runs outside the capture callback. The streaming receive worker validates
finite samples before advancing state, bounds processing steps to 480 samples,
and retains at most 30 s of mono history (1,440,000 float samples, about
5.76 MB of sample payload). Demodulation allocates additional bounded working
arrays, so the retained-history limit is not a whole-process memory limit.
Absolute half-open sample ranges survive chunk boundaries and ring-buffer
wrap. If a candidate began before retained history, it is marked truncated and
emits no partial frame. WAV reports count truncated energy candidates and
retain independently acquired pilot frames.

The energy detector accumulates 10 ms windows. A window is active when
\[
 \frac{1}{W}\sum_{n=0}^{W-1}x[n]^2 \geq 10^{-4},\qquad W=480.
\]
A 10 ms window holds 480 samples at 48 kHz. Square each normalized
sample so positive and negative excursions both count, then average.
A mean square of $10^{-4}$ corresponds to an RMS amplitude of 0.01,
so this gate opens at one percent of full scale.

A quiet window closes a candidate; end-of-stream flushes a partial window and
any open candidate. The 1\% full-scale RMS gate serves level-4 acquisition
with an external measurement profile and reports energy candidates for
diagnostics. The CLI scans level-1--3 pilots independently of this gate.
Energy candidates can include impulsive noise and unrelated signals.

\subsection{BW500 levels 1--3: tone-index geometry}
\modemcite{proakis}These levels encode data by selecting a frequency bin for each symbol. They
are not the level-4 differential-BPSK waveform played more slowly. The
implementation performs direct sine/cosine dot products at selected DFT bins;
it does not compute a full FFT for each symbol.

\begin{table}[htbp]
\centering\small
\begin{tabular}{@{}lrrrrrrr@{}}
\toprule
Level & $H$ & Bins & Columns & Refs. & $N$ & Coded & Payload\\
\midrule
1 & 2048 & 54--74 & 124 & 24 & 96 & 300 & 9 B\\
2 & 1024 & 27--37 & 228 & 24 & 200 & 612 & 22 B\\
3 & 1024 & 27--37 & 226 & 24 & 297 & 606 & 34 B\\
\bottomrule
\end{tabular}
\caption{Implemented index-mode geometry at 48 kHz. Columns exclude the initial lead; $N$ includes mode padding before whitening. Payload excludes the control byte and two-byte CRC.}
\label{tab:dsp-index-geometry}
\end{table}

The symbol rates are $48000/H$: 23.4375 symbols/s at level 1 and 46.875 at
levels 2/3. The synthetic transmitter adds two lead symbols at level 1 and four
at levels 2/3, all totaling 4,096 samples. Those bench leads use fixed in-band
tones; they do not implement the stock session's randomized lead stream.
There are three coded bits per data column. Subtracting the 24 reference
columns gives $100\times3=300$, $204\times3=612$, and $202\times3=606$
channel bits respectively.

Mode construction uses seeds 12,600,710; 7,910,701; and 14,955,604. One generator
stream produces the channel permutation, then the turbo permutation, followed
by bin/pilot draws. Level 1 replaces its generated base bins with the explicit
measured 124-entry array in the source. Levels 2/3 currently use generated
base allocations. All three therefore need qualification appropriate to their
actual construction, not a blanket claim that every allocation is independently
measured.

For each data column, the eight labels $(0,1,3,2,7,6,4,5)$ select offsets from
the column's base bin. Level 1 applies stride 2 modulo a span of 21; levels
2/3 apply stride 1 modulo 11. Fixed reference columns and their known bin
values allow synchronization without comparing decoded payload to an expected
message. The level-3 information block contains 37 complete bytes plus one
padding bit: 34 payload bytes, one control byte, two CRC bytes, and a zero
padding bit before whitening.

\landscapeplantuml{build/diagrams/Linkwise_DSP_Index}{Implemented BW500 level 1--3 receive pipeline. Pilot acquisition locates index frames in continuous audio. Soft demodulation and CRC validation follow timing and frequency selection.}{fig:dsp-index}

\subsection{Index-mode acquisition and soft demodulation}
\modemcite{proakis,meyr}The WAV and device receive commands use \path{modem::bw500::scan} to
search every 128 samples in a bounded rolling spectrum history. At each
epoch, the scanner scores the 24 reference columns against energy in
the occupied band and guard bins at offsets of 0 and $\pm$12 Hz. It
considers hypotheses with a mean pilot score of at least 0.65 and
refines timing in 16-sample steps. This lets a frame span brief drops in
energy or continuous receiver noise. The library also retains the
candidate-based decoder, which starts its search near $b+LH$ for burst
start $b$ and lead length $L$. The score threshold has no measured
false-alarm rate.

For a symbol of length $H$ and bin $k$, the unnormalized tone energy is
\[
 E_k=\left|\sum_{n=0}^{H-1}x[n]\exp(-j2\pi kn/H)\right|^2.
\]
The exponential supplies a reference that completes $k$ cycles over the
$H$ samples. Multiplying and summing aligns a matching tone so its
contributions add together. Taking the squared magnitude measures the
strength of that match independently of its starting phase. For example,
at 48 kHz with $H=1024$, bin 32 tests for a 1,500 Hz tone.

For timing refinement and soft decisions, the implementation divides by
$(H/2)\sum x[n]^2$, with an epsilon floor. For data, it measures the
eight mapped bin hypotheses. For each of the three bit positions, the soft
metric is the largest normalized power among labels with bit 1 minus the
largest among labels with bit 0. This max-based difference is a decoder input;
it is not a calibrated log likelihood, SNR, or BER estimate.
A positive value means the strongest tone compatible with bit 1 beat
the strongest tone compatible with bit 0. Values near zero give the
decoder little reason to prefer either bit.

Metrics are scattered through the inverse channel permutation into FEC order.
Only a CRC-valid result is returned. The rolling scanner suppresses overlap
with an emitted frame and continues to search for adjacent frames. Level 4
handles one frame per energy candidate. Continuous carrier and sample-clock
tracking remain open.

\subsection{Index-mode transmit synthesis}
The level-1/2/3 transmitter accepts a nonempty list of frames with exactly
9, 22, or 34 payload bytes per frame. It appends the control byte and byte CRC,
serializes MSB-first, pads to the mode's $N$, whitens, and turbo-encodes. For
each data column it gathers three channel-ordered coded bits into an integer
0--7 and selects the corresponding mapped tone. Reference columns emit their
fixed expected bins without consuming coded bits. A burst has one lead at its
start, followed by consecutive frame-column blocks; the lead is not repeated
between frames.

Each selected bin $k$ is rendered as $0.5\sin(2\pi kn/H)$ for $n=0,\ldots,H-1$.
The phase argument restarts for each symbol; integer-bin tones complete an
integer number of cycles over $H$ samples. The implementation concatenates
those symbol buffers directly. It does not add DATA-symbol pulse shaping,
resampling, a stock randomized acquisition stream, or a qualified transmitter
spectral mask. The fixed half-scale amplitude is a bench sample convention.
Reference-peer acceptance, occupied bandwidth, audio-route gain, and actual
RF power must be measured separately before treating this synthesizer as an
operational transmitter.

\subsection{Shared CRC, whitening, and interleaver generation}
\modemcite{koopman}Frame integrity uses CRC-16/GENIBUS: polynomial \texttt{0x1021}, initial register
\texttt{0xffff}, final complement, no reflection, and a big-endian two-byte
trailer. The check vector is \texttt{123456789} $\mapsto$ \texttt{0xd64e}.
The CRC covers the actual frame body, including its control field; mode padding
bits are outside the byte CRC. CRC checks establish frame integrity under this
code, not authenticated sender identity or recipient delivery.

Bits are serialized most-significant first, padded to the mode's information
size, and XORed with a freshly seeded whitening sequence. The 24-bit generator
obeys
\[
 s_{i+1}=(\texttt{0x43fd43fd}\,s_i+\texttt{0xc39ec3})\bmod 2^{24}.
\]
Each step multiplies the current state $s_i$ by the first constant,
adds the second constant, and keeps the lowest 24 bits. The remainder
operation is what makes the state wrap. The same starting state and
steps produce the same bit sequence at both ends of the link.

Whitening starts at $s_0=29318$ and takes the next state's high bit. A
generator bit of 1 flips the corresponding frame bit. A generator bit
of 0 leaves it alone. Applying the same mask again flips those bits
back. Permutations use
rejection sampling: draw an index, append only its first occurrence, and
continue until every index is present. The draw multiplies the normalized
state by the requested size with single-precision rounding, matching the
implementation's wire-oriented generator semantics. Changing numeric precision
or substituting a library shuffle can change the permutation.

The channel permutation and turbo permutation are different objects. If
$\pi_1$ maps transmitted position $j$ to a coded-bit position, TX gathers
$q_j=c_{\pi_1(j)}$ and RX scatters the corresponding soft metric back into
$c_{\pi_1(j)}$. The inner permutation $\pi_2$ changes the input order of the
second constituent encoder and must be inverted consistently when exchanging
extrinsic information. Interleavers belong to the mode/profile, not to the
agent's request parameters.
In the worked example, the channel permutation puts coded bit 33 in the
first transmitted position. The receiver puts the first soft decision
back into slot 33 before decoding. The turbo permutation acts earlier,
when the second encoder reads the information bits in a different order.

\subsection{Turbo encoder and iterative receive decoder}
\modemcite{berrou}The codec has two eight-state recursive systematic convolutional constituents.
Writing the three state bits as $(r_0,r_1,r_2)$ and the next information bit as
$u$, the implementation uses
\[
 a=u\oplus r_1\oplus r_2,\quad
 p=a\oplus r_0\oplus r_2,\quad
 (r_0',r_1',r_2')=(a,r_0,r_1).
\]
The symbol $\oplus$ means exclusive OR. First compute the new feedback
bit $a$ from the input and two stored bits. Then compute the parity bit
$p$ from $a$ and two stored bits. Finally, shift the memory so $a$
becomes the newest bit. Starting with an empty register and input 1,
the encoder outputs parity 1 and moves to state $(1,0,0)$.

Both encoders start at zero. The first consumes natural-order bits and the
second consumes turbo-permuted bits. Levels 1/2 transmit systematic and both
parity streams (rate 1/3); levels 3/4 transmit systematic bits and alternate
which constituent parity is retained (rate 1/2). Each constituent is terminated
with three steps using $u=r_1\oplus r_2$. The wire tail is twelve bits,
interleaving the two systematic/parity tail pairs at each of three steps.
Thus the implemented coded sizes are $3N+12$ or $2N+12$.
The factor of three or two counts the transmitted bits per information
bit before termination. For the worked level-4 frame, 368 information
bits give 736 systematic and parity bits, then the tail brings the
total to 748.

\landscapeplantuml{build/diagrams/Linkwise_DSP_FEC}{Implemented shared whitening, constituent coding, tail and permutation structure. Receive decoding uses a hard-decision CRC fast path followed by at most eight max-log-MAP iterations.}{fig:dsp-fec}

\modemcite{bcjr,robertson}The receiver first makes hard decisions on systematic metrics and checks the
frame CRC. If that fails, it constructs the two terminated trellises, treating
punctured parity as a zero-information metric. A forward/backward
\emph{max-log-MAP} pass computes extrinsic information for each constituent.
The implementation uses maxima rather than log-sum-exp correction; it must
not be described as an exact log-MAP decoder. Positive input metrics favor
bit 1. The two constituents exchange interleaved/de-interleaved extrinsic
information for at most eight iterations, checking de-whitened frame integrity
after each iteration and stopping on a valid CRC.

This implementation does not maintain cross-retransmission soft combining.
Its metrics have heuristic scales and its tests do not establish coding gain,
waterfall curves, or residual undetected-error probability under HF fading.
Those require independently generated vectors and controlled channel sweeps,
including offset, burst noise, clipping, and long erasures.

\subsection{Level-4 frame geometry and external measurement profile}
The level-4 byte container is 46 bytes: 43 payload bytes, one control byte,
and two CRC bytes. Its $N=368$ information bits produce 748 coded bits. The
channel and turbo permutations start from seed 6,437,280. Pilot generation
continues that stream after the two permutations, choosing 18 distinct columns
per band with $5+\lfloor R\,386\rfloor$.

The implemented waveform stores 393 columns per band: one reference followed
by 392 body columns. Removing 18 pilot positions in each band yields
$2(392-18)=748$ differential data cells. This is the retained single-frame
implementation geometry, not a claim that every stock burst has exactly that
many emission slots. Its reference centers start after nine 512-sample slots.
Band centers are 1,350 and 1,650 Hz, the symbol hop is $H=512$ samples
(93.75 symbols/s), and the upper band is staggered by 256 samples. Unlike the
index ladder, these carriers are not selected integer bins of a symbol FFT.

The phase reference consists of 786 measured points on a 128-phase lattice,
393 per band. Each point is $G_{b,t}=\exp(j2\pi k_{b,t}/128)$. A shared measured
513-sample real pulse covers offsets $-256$ through $+256$. The profile loader
requires all reference positions, rejects duplicates and out-of-range phases,
checks pulse length and finite/bounded taps, and limits JSON size to 128 KiB.
Those structural checks do not prove that an arbitrary supplied profile matches
a real modem. Reproducible tests use SHA-256-pinned upstream measurements,
extracted without executing downloaded source. Cached measurement sources
retain their own attribution/license metadata and are not embedded in Rust.

The phase index $k_{b,t}$ selects one of 128 equally spaced directions
around a circle. For example, index 32 represents a quarter turn. The
transmitter and receiver use opposite rotations so they can compare
successive DATA symbols after accounting for the known phase pattern.

\landscapediagram{\begin{tikzpicture}[x=1.02cm,y=1cm,font=\small,
  axis/.style={ovgray,-{Stealth[length=2mm]}},
  data/.style={ovblue,fill=ovblue},
  pilot/.style={ovgold!60!black,fill=ovgold},
  reference/.style={ovteal!70!black,fill=ovteal}]
  \node[anchor=west,font=\bfseries] at (-2.7,6.15) {BW500 level-4 symbol centers};
  \node[anchor=west] at (-2.7,5.55) {$n_{b,j}=4608+512j+256b$, with band index $b\in\{0,1\}$ and column $0\le j\le392$};
  \foreach \x in {0,1,...,17}
    \draw[ovgray!20] (\x,1.05) -- (\x,4.95);

  \draw[axis] (0,3.9) -- (17.65,3.9);
  \draw[axis] (0,2.45) -- (17.65,2.45);
  \node[align=right,anchor=east] at (-0.35,3.9) {Lower band\\1,350 Hz};
  \node[align=right,anchor=east] at (-0.35,2.45) {Upper band\\1,650 Hz};
  \draw[reference] (-0.09,3.81) rectangle ++(0.18,0.18);
  \draw[reference] (0.91,2.36) rectangle ++(0.18,0.18);
  \foreach \j in {1,...,8} {
    \fill[data] ({2*\j},3.9) circle[radius=2.4pt];
    \ifnum\j=7
      \draw[pilot] (15,2.59) -- (15.14,2.45) -- (15,2.31) -- (14.86,2.45) -- cycle;
    \else
      \fill[data] ({2*\j+1},2.45) circle[radius=2.4pt];
    \fi
  }
  \node[above,font=\scriptsize] at (0,4.11) {$j=0$};
  \node[above,font=\scriptsize] at (2,4.11) {$j=1$};
  \node[above,font=\scriptsize] at (15,2.59) {$j=7$, first upper pilot};
  \draw[ovblue,{Stealth[length=1.4mm]}-{Stealth[length=1.4mm]}]
    (0,4.7) -- (2,4.7) node[midway,above,font=\scriptsize] {512-sample hop};
  \draw[ovteal!70!black,{Stealth[length=1.4mm]}-{Stealth[length=1.4mm]}]
    (0,3.15) -- (1,3.15) node[midway,above,font=\scriptsize] {256};

  \draw[axis] (0,1.05) -- (17.65,1.05);
  \foreach \x/\sample in {0/0,2/512,4/1024,6/1536,8/2048,10/2560,12/3072,14/3584,16/4096}
    \draw[ovgray] (\x,1.05) -- ++(0,-0.11) node[below,font=\scriptsize] {\sample};
  \node[anchor=west] at (0,0.34) {Samples relative to the lower-band reference, $n-4608$};

  \draw[reference] (0,-0.55) rectangle ++(0.18,0.18);
  \node[anchor=west] at (0.32,-0.46) {Reference};
  \fill[data] (3.05,-0.46) circle[radius=2.4pt];
  \node[anchor=west] at (3.27,-0.46) {Data};
  \draw[pilot] (5.1,-0.32) -- (5.24,-0.46) -- (5.1,-0.60) -- (4.96,-0.46) -- cycle;
  \node[anchor=west] at (5.34,-0.46) {Pilot};
  \node[anchor=west,font=\scriptsize] at (7.0,-0.46)
    {18 pilots per band. The first lower-band pilot is at $j=24$.};
\end{tikzpicture}
}{BW500 level-4 symbol centers on the 48 kHz sample axis. The plotted excerpt uses the implemented pilot allocation. Each marker locates the center of a 513-tap pulse extending 256 samples on either side.}{fig:dsp-level4-symbol-timing}

\subsection{Level-4 transmit synthesis}
After whitening, turbo encoding, and channel gather, coded bits fill band 0's
data positions in time order, then band 1's. With initial polarity
$v_{b,0}=+1$, a coded 1 negates polarity, a coded 0 preserves it, and a pilot
repeats the preceding polarity. Each complex emitted cell is
$c_{b,t}=v_{b,t}\overline{G_{b,t}}$. The real audio is a sum of pulse-shaped
carriers:
\[
 x[n]=\frac14\sum_{b=0}^{1}\sum_t p[n-n_{b,t}]
       \Re\{c_{b,t}\exp(j2\pi f_b n/48000)\},
\]
where $n_{b,t}=9H+tH+bH/2$, $f_b\in\{1350,1650\}$, and $p$ has 513 taps.
To compute one output sample $x[n]$, center each symbol's pulse on
$n_{b,t}$, evaluate the carrier at sample $n$, and add the real
contributions from both bands. The pulse determines how much each nearby
symbol contributes. The center formula allows nine hops before the
first reference, advances one hop per column, and offsets the upper
band by half a hop.

The factor $1/4$ is the current bench amplitude scale, not a calibrated radio
power command. A waveform's audio amplitude does not replace CAT power limits,
route qualification, or PTT authority.

The synthetic buffer contains $9H+393H+1=205825$ samples, approximately
4.288 s at 48 kHz. The initial slots do not contain a generated randomized
acquisition lead. This output length is a local implementation fact and must
not be substituted for a measured reference modem's PTT interval, framing
duration, or turnaround requirement. Actual device playback and RF keying are
not performed by the file-only synthesis tool.

\subsection{Level-4 timing lock, demodulation, and validation}
\modemcite{proakis,meyr}The receiver mixes the real input down at each nominal carrier using
$x[n]\exp(-j2\pi f_b n/48000)$. At a trial symbol center it multiplies the
513 surrounding complex samples by the measured real pulse and sums them.
It then multiplies by $G_{b,t}$ to remove the known phase reference. If the
result is $z_{b,t}$, the differential statistic is
\[
 d_{b,t}=\Re\{z_{b,t}\overline{z_{b,t-1}}\}.
\]
This dot product compares consecutive symbols after removing the known
reference phase. It tends to be positive when their polarities agree
and negative when they are opposite. A DATA bit 1 flips polarity,
so the receiver negates this statistic to give a positive score for 1.

A correct pilot transition should preserve polarity. Timing search maximizes
the average normalized differential dot product over 36 pilot positions,
first in 32-sample steps and then in two-sample steps around the best epoch.
A score below 0.65 is rejected. The search is bounded by the candidate, frame
geometry, available samples, and one-hop end tolerance; it does not scan an
unbounded recording for a matching payload.

Data positions produce $-d_{b,t}$ so positive metrics favor a transmitted 1.
The 748 values are scaled by six divided by their mean absolute magnitude,
then passed through channel scatter and the shared turbo/de-whitening/CRC
chain. This normalization is not an estimate of a physical noise variance.
No valid CRC means no frame. A valid frame is interpreted in protocol context:
a setup frame can identify a caller; a DATA container can carry payload and
control; neither is an unconditional link-up event.

\landscapeplantuml{build/diagrams/Linkwise_DSP_Level4}{Implemented level-4 transmit and receive transforms. Measured reference phases and pulse taps are external profile inputs; the single-frame synthetic round trip does not qualify reference interoperability.}{fig:dsp-level4}

The level-4 decoder currently allocates complex mixed buffers and performs
repeated direct matched-pulse sums. It is a bench implementation, not a
profiled real-time DSP scheduler. An optimized runtime may cache oscillators,
reference terms, pulse weights, and mode tables, and bound search work per
candidate, but must preserve sample epochs, metric sign, permutation order,
and CRC outcomes against golden fixtures before replacing these kernels.
